\subsection{Four-state act cones through corner order fourteen}

\begin{theorem}[Proved: four states and fourteen corner elements]
\label{CP2:A-four-state-fourteen}
Let $M$ be a finite monoid of order at most fourteen and $S$ a
finite right $M$-act with at most four states. The ordinary
two-object act-cone category $\mathcal C(M,S)$ cannot have
classifying space of type $\CP^2$. The left-act assertion holds
by opposition. In particular a four-state ordinary act cone
of type $\CP^2$ has at least twenty arrows, identities counted.
\end{theorem}
\begin{proof}
For at most three states use Theorem~\ref{CP2:A-three-cross-reduction}
and the three-state minimum; for four states and $|M|\leq12$
use Corollary~\ref{CP2:A-four-state-small}. It remains to consider
$|M|=m\in\{13,14\}$ and $|S|=4$.

Suppose the cone has projective-plane type. The relative-cup
obstruction requires nonzero rational second cohomology of $BM$,
and the two-row act obstruction requires at least three minimal
rows. Its minimal ideal therefore has $r\geq3$ rows and $c\geq2$
columns. If its maximal group has order at least two, the ideal
has at least twelve elements, leaving at most two outside
elements. Lemma~\ref{CP2:A-two-outside-acts} then gives zero
fourth rational homology from boundedness, a contradiction.
Thus the minimal ideal is the rectangular band $I\times J$,
of shape $r\times c$. Write $q=m-rc\geq1$ for the number of
outside elements; the identity is outside this proper ideal.

Lemma~\ref{CP2:A-four-state-band} says that either every rational
degree-two product is zero or the band acts constantly. The
first alternative is impossible. In the second, the actual
compression of Theorem~\ref{CP2:A-constant-act-compression}
gives a monoid $N$ of the same homotopy type, with minimal band
$I\times S$, the same $q$ outside elements, and
\[
 |N|=q+4r=m+r(4-c).
\]
If $c\geq3$, this is at most $m+\lfloor(m-1)/3\rfloor\leq18$
when $c=3$, and at most $m$ when $c\geq4$. The complete monoid
minimum excludes both. We may therefore assume $c=2$.

For $r=3$ or $4$, the compressed order is at most $22$, again
excluded. For $r\geq6$, the original outside count is at most
$14-12=2$, and the two-outside lemma applies directly to the
original act. Only $r=5$ remains. When $m=13$, the compressed
monoid has order $23$ and minimal band $5\times4$. This
contradicts the complete classification
Theorem~\ref{A:monoid-minimum23-classification}: the minimal-band
shapes of the two ordinary minimizing monoids are $3\times5$
and $5\times3$.

Finally let $m=14$, $r=5$, $c=2$, $q=4$. For the compressed
monoid put $A=\mathbb QN/\mathbb Q(I\times S)$. This is the
same contracted outside algebra as that of $M$, of dimension
four. Its row and column augmentation modules, respectively,
are the left $A$-module $V$ of dimension four and the right
$A$-module $W$ of dimension three. The band--Tor formula
would require
\[
 \dim(W\otimes_A V)=1,\qquad
 \dim\operatorname{Tor}_2^A(W,V)=1,\qquad
 \operatorname{Tor}_j^A(W,V)=0\quad(j\geq3).
\]
If $A$ is local, Lemma~\ref{A:local-target-facts} makes $V$
cyclic. Its cyclic cover $A\to V$ has equal source and target
dimensions and is therefore an isomorphism. This makes $V$
projective and contradicts the required second Tor. If $A$
is nonlocal, Lemma~\ref{A:nonlocal-four} says that nonzero
second Tor persists in arbitrarily high degrees, again a
contradiction. This exhausts the two remaining monoid orders.
Since a four-state ordinary act cone has $m+5$ arrows, the
last assertion follows. Every compression used above is an
actual homotopy equivalence; no rational comparison was
promoted to a homotopy equivalence.
\end{proof}
